% !PRINTED 51
\weberpage{51}

%<<CONT
and applied once again:
\weq{9}{\begin{gathered}
\mathrm{sn}\,(u+4K) = \mathrm{sn}\,u,\\
\mathrm{cn}\,(n+4K) = \mathrm{cn}\,u, % [sic] \textquotedbl{}n\textquotedbl{} printed for \textquotedbl{}u\textquotedbl{}
\\
\mathrm{dn}\,(u+4K) = \mathrm{dn}\,u.
\end{gathered}}
The elliptic functions therefore have the property of \gesperrt{periodicity} in common with the trigonometric functions. They have the period $4K$ (dn $u$ also the period $2K$).

In §\,3 we derived, as an example, the transformation, formula (1):
\weq{10}{\frac{dz}{\sqrt{z(1-z)(1-\varkappa^2 z)}} = \frac{dx}{i\sqrt{x(1-x)(1-\varkappa'^2 x)}},}
when
\weq{11}{z = \frac{-x}{1-x}}
held. If we therefore set
\[
\frac{dx}{\sqrt{x(1-x)(1-\varkappa'^2x)}} = i\,du,
\]
then, if $x$ and $u$ vanish together, $x = \mathrm{sn}^2(iu,\varkappa')$, $1-x = \mathrm{cn}^2(iu,\varkappa')$, and from (10) it follows that
\[
\frac{dz}{\sqrt{z(1-z)(1-\varkappa^2z)}} = du,
\]
and hence
\[
z = \mathrm{sn}^2(u,\varkappa).
\]
Accordingly (11) gives
\weq{12}{\mathrm{sn}\,(u,\varkappa) = \frac{-i\,\mathrm{sn}\,(iu,\varkappa')}{\mathrm{cn}\,(iu,\varkappa')}}
(in which the sign is to be determined from the special value $u=0$).

If we now set
\[
\int_0^{\frac{\pi}{2}} \frac{d\varphi}{\sqrt{1-\varkappa'^2\sin^2\varphi}} = K',
\]
then from (8) there results
\[
\begin{aligned}
\mathrm{sn}\,(iu+2K',\varkappa') &= -\mathrm{sn}\,(iu,\varkappa'),\\
\mathrm{cn}\,(iu+2K',\varkappa') &= -\mathrm{cn}\,(iu,\varkappa'),
\end{aligned}
\]
and consequently from (12)
\[
\mathrm{sn}\,(u-2iK',\varkappa) = \mathrm{sn}\,(u,\varkappa)
\]
%CONT>>
